\def\stat{simakov}



\def\tit{МНОГОКРИТЕРИАЛЬНЫЙ МЕТОД ВЫЯВЛЕНИЯ НЕЧЕТКИХ ДУБЛИКАТОВ

В~ПОТОКЕ ТЕКСТОВЫХ СООБЩЕНИЙ$^*$}



\def\titkol{Многокритериальный метод выявления нечетких дубликатов

в~потоке %текстовых

сообщений}



\def\autkol{А.\,М.~Андреев, Д.\,В.~Березкин, И.\,А.~Козлов,

К.\,В.~Симаков}



\def\aut{А.\,М.~Андреев$^1$, Д.\,В.~Березкин$^2$, И.\,А.~Козлов$^3$,

К.\,В.~Симаков$^4$}



\titel{\tit}{\aut}{\autkol}{\titkol}



{\renewcommand{\thefootnote}{\fnsymbol{footnote}}

\footnotetext[1]

{Статья рекомендована к публикации в~журнале Программным комитетом конференции <<Электронные

библиотеки: перспективные методы и~технологии, электронные коллекции>> (``Digital Libraries: Advanced

Methods and Technologies,'' RCDL-2014), г.~Дубна, 13--16~октября 2014~г.}}



\renewcommand{\thefootnote}{\arabic{footnote}}

\footnotetext[1]{Московский государственный технический университет им.\ Н.\,Э.~Баумана,\\ arkandreev@gmail.com}

\footnotetext[2]{Московский государственный технический университет им.\ Н.\,Э.~Баумана,\\ dmitryb2007@yandex.ru}

\footnotetext[3]{Московский государственный технический университет им.\ Н.\,Э.~Баумана,\\ kozlovilya89@gmail.com}

\footnotetext[4]{Московский государственный технический университет им.\ Н.\,Э.~Баумана, skv@ixlab.ru}





\label{st\stat}



                  \thispagestyle{headings}



\Abst{Рассмотрена задача обнаружения нечетких дубликатов в~потоке

текстовых сообщений. Предложена модель документа, имеющая

возможность гибкой настройки на различные предметные области.

Представлен многокритериальный метод выявления дублирующихся

документов на основе бинарной классификации с~помощью метода опорных

векторов. Предложен способ обеспечения высокого быстродействия метода

посредством предварительного отбора кандидатов в~дубликаты. Проведена

экспериментальная оценка предложенного метода, демонстрирующая его

практическую применимость.}



\KW{обнаружение нечетких дубликатов; мера близости; бинарная

классификация}



\DOI{10.14357\!/\!08696527150103}







\vskip 12pt plus 9pt minus 6pt





      \thispagestyle{headings}



\section{Введение}



  В настоящее время во многих предметных областях существует потребность

в формировании больших текстовых коллекций. При этом производится сбор

текстовой информации из открытых ин\-тер\-нет-ис\-точ\-ни\-ков, а~также

специализированных ресурсов. Основной областью использования создаваемых

таким образом хранилищ документов является интеллектуальная обработка

текстов, которую, как правило, можно отнести к классу Text Mining.



  С ростом количества разнообразных источников данных в~сети Интернет

(новостные сайты, блоги, социальные сети) все более серьезной проблемой

становится дублирование информации. Сообщения, публикуемые одним

источником, зачастую многократно перепечатываются другими (в исходном

виде или с~небольшими изменениями). В~результате при выполнении

автоматического сбора документов из многочисленных источников в~формируемой текстовой коллекции накапливаются идентичные или близкие по

содержанию документы~--- дубликаты. В~некоторых задачах наличие

дубликатов должно учитываться~--- например, при определении значимости

сообщений~[1]. Но в~большинстве случаев попадание таких документов в~коллекцию снижает ее качество~[2].



  В данной статье рассматривается решение задачи обнаружения дубликатов в~потоке текстовых сообщений. Особое внимание уделяется обеспечению

возможности использования разработанного метода для обработки документов

из различных предметных областей.



\section{Постановка задачи}



\subsection{Функционирование системы сбора и~обработки новостной

информации}



  В работе~[3] авторами было предложено решение задачи качественного

автоматического сбора новостных данных из интернет-источников,

предполагающего извлечение с~веб-стра\-ни\-цы текста новости, а~также

сопутствующих метаданных, включающих название, дату публикации, автора

новости и~др. Текстовые данные, извлеченные с~веб-сай\-тов, подвергаются

последующей обработке различными методами интеллектуального анализа,

такими как автоматическая классификация и~кластеризация документов,

извлечение знаний и~фактов из ес\-те\-ст\-вен\-но-язы\-ко\-вых текстов и~прогноз

развития ситуаций.



  Для эффективного применения этих методов необходимо обеспечить

качество анализируемой коллекции документов, в~частности~--- устранить из

нее дубликаты. Поэтому в~систему сбора требуется встроить подсистему,

выполняющую их оперативное обнаружение и~удаление (рис.~1). Она должна

сравнивать каждый загружаемый документ с~полученными ранее новостями и~удалять его, если он является дубликатом одной из них.







\subsection{Особенности решаемой задачи}



  Если проблема обнаружения и~удаления полных дублей тривиальна, то при

необходимости распознавать нечеткие дубликаты (т.\,е.\ документы, имеющие

различный текст, но близкие по содержанию) возникают значительные

сложности.



  В связи с~явлением частичного дублирования в~работе~[4] предлагается

понятие информативной необходимости элемента высказывания: <<Всякий

час\-тич\-но-дуб\-лет\-ный элемент, если он отвечает коммуникативной задаче

высказывания, признается информативно-необходимым в~той же мере, что и~прочие, не\-дуб\-лет\-ные элементы речи: такой элемент информативно необходим

постольку, поскольку вносит свой уникальный, неповторимый в~других

элементах, вклад в~суммарную информацию, передаваемую высказыванием>>.

Таким образом,  задача обнаружения нечетких дубликатов состоит в~распознавании и~удалении сообщений, не являющихся

  информативно-необходимыми.







  Установить информативную необходимость и~ценность некоторого

высказывания можно только с~учетом соответствующего ситуативного

контекста. Так, при ручной проверке документов эксперт, определяя наличие

или отсутствие дублирования, принимает решение с~учетом предметной

области и~характера документов, составляющих анализируемую коллекцию.

    \begin{figure} %fig1

        \vspace*{1pt}

 \begin{center}

 \mbox{%

 \epsfxsize=117.482mm

 \epsfbox{sim-1.eps}

 }

 \end{center}

 \vspace*{-10pt}

  \Caption{Место подсистемы обнаружения дубликатов в~системе автоматизированного

сбора и~анализа новостной информации}

\vspace*{-9pt}

  \end{figure}

  Например, при работе с~юридическими документами особое внимание должно

уделяться метаданным~--- для текстов такого типа два документа с~практически

идентичным содержанием, но различающимися названиями и~датами

публикации не могут считаться дублями. Другой подход используется при

анализе экономических новостей, например сводок о состоянии фондового

рынка,~--- дубликатами не должны признаваться документы, имеющие одно и~то же название и~одинаковый текст, но различающиеся числовыми данными

(значениями курсов валют).



  Таким образом, в~разных случаях эксперт сравнивает документы с~точки

зрения различных критериев (близость текстового содержания, сходство

названий, разница во времени публикации), т.\,е.\ использует различные модели

распознавания дубликатов в~зависимости от предметной области и~контекста

коммуникационных сообщений. Поэтому не представляется возможным

выработать единую систему правил для обнаружения нечеткого дублирования

сразу для всех случаев. Следовательно, разрабатываемая подсистема должна

иметь возможность гибкой настройки на разные предметные области, что

позволит ей моделировать деятельность эксперта по распознаванию дублей с~использованием различных моделей. Поскольку система сбора выполняет

извлечение документов из множества источников, которые относятся к

различным предметным областям, необходимо обеспечить возможность работы

с несколькими моделями распознавания дубликатов одновременно.



  Еще одной проблемой, с~которой приходится сталкиваться при решении

задачи устранения дубликатов, является большой объем обрабатываемых

данных. Наличие в~коллекции сотен тысяч документов делает весьма

трудоемким анализ каждого нового сообщения путем сравнения его с~каждым

из ранее загруженных. Эта проблема может быть решена с~помощью

приближенных методов, но их применение ведет к снижению качества

обнаружения дубликатов, т.\,е.\ уменьшению точности и~полноты~[2]. При

разработке предлагаемого подхода решалась задача совмещения высокого

качества и~низкой вычислительной сложности проверки документов.



\vspace*{-4pt}



\section{Обзор методов обнаружения дублей}



\vspace*{-4pt}



\subsection{Однокритериальные методы}



\vspace*{-2pt}



  Большинство методов, используемых при выявлении дублирования

сообщений, основано на общей идее~--- попарном сопоставлении документов с~использованием выбранной меры близости. При этом документы~$d_i$ и~$d_j$

считаются дубликатами, если значение этой меры для них превосходит

заданную пороговую величину: $\mathrm{sim}\left(d_i,d_j\right)\hm>\mathrm{Tr}$. Методы, предлагаемые

исследователями, различаются способом представления документов и~используемыми мерами близости.



  Наиболее эффективным с~точки зрения вычислительной сложности являются

представление документа одним числовым значением~---  <<сигнатурой>>~---

и~определение дублирования сообщений на основании совпадения их сигнатур.

В~качестве сигнатур используются:

  \begin{itemize}

\item контрольная сумма строки, полученной из нескольких сцепленных в~алфавитном порядке слов документа с~наибольшими значениями весов (для

расчета весов используются методы TF (term frequency),

TF-IDF (inverse document frequency) и~OptFreq (optimal

frequency))~[2];

\item контрольная сумма строки, полученной из нескольких наиболее длинных или

значимых (т.\,е.\ состоящих из слов с~наибольшим суммарным значением весов)

предложений документа~[2];

\item контрольная сумма строки, полученной из слов текста, принадлежащих

заданному словарю (содержащему слова со средними значениями IDF)~[5];

\item битовый вектор, каждый элемент которого соответствует слову из заданного

словаря и~равен~1, если частота встречаемости этого слова в~документе

превышает пороговое значение, и~0~в~противном случае~[6].

  \end{itemize}



  Другой популярный подход состоит в~представлении документа множеством

элементов. Близость множеств, соответствующих двум сообщениям,

рассчитывается, как правило, с~помощью коэффициента Жаккара. В~качестве

элементов множеств используются:

  \begin{itemize}

\item предложения документа~[7];

\item ключевые термы (слова, наиболее значимые с~точки зрения

лингвостатистических критериев)~[1];

\item шинглы (всевозможные последовательности фиксированной длины~$k$,

состоящие из соседних слов документа)~[8];

\item N-граммы (всевозможные последовательности, состоящие из $N$ соседних

символов документа)~[9].

  \end{itemize}



  Широко применяются векторные модели документов, в~которых каждое

сообщение представляется вектором в~многомерном признаковом

пространстве. В~качестве элементов вектора обычно используются слова,

встречающиеся в~текстах коллекции~[10]. При этом значениями элементов

являются веса соответствующих слов, отражающие их значимость для

документа. В~качестве признаков могут применяться и~другие характеристики

текстов, например частота появления различных пар символов или частота

появления тех или иных частей речи~[11]. Для сравнения векторов часто

используется косинусная мера или ее модификации~[10].



  Существует несколько методов, в~которых элементами вектора являются

сигнатуры, полученные различными способами. При этом для сравнения

векторов используется расстояние Хэмминга. Так, в~\cite{13-sim} предложена

модификация сигнатурного метода~\cite{5-sim}, предполагающая использование

нескольких словарей, для каждого из которых вычисляется своя сигнатура.

Другим примером этого подхода служит метод, развивающий идею алгоритма

шинглов~\cite{8-sim, 14-sim}. Он состоит в~применении к элементам множества

шинглов различных хэш-функ\-ций и~выборе для каждой из них шингла,

минимизирующего ее значение. Из выбранных шинглов формируются группы,

именуемые <<супершинглами>>, которые и~являются элементами вектора,

представляющего документ.



  Наконец, в~\cite{12-sim} предлагается использовать в~качестве модели

документа упорядоченную последовательность символов. Для сравнения в~этом

случае используются редакционные расстояния, такие как расстояние

Левенштейна.



\subsection{Многокритериальные методы}



  В отличие от вышеперечисленных подходов методы, описанные

  в~\cite{12-sim, 15-sim}, основаны на оценке схожести документов с~точки

зрения нескольких критериев одновременно. Это связано с~тем, что

анализируемые данные в~этих работах состоят из множества полей и~для

каждого поля значение близости записей рассчитывается отдельно (при этом

для разных полей могут использоваться различные меры близости).

В~результате пара записей представляется вектором, каждый элемент которого

соответствует схожести соответствующих полей документов. Пары,

отмеченные в~обучающей выборке как <<дубликаты>> или <<недубликаты>>,

представлены двумя кластерами точек многомерного пространства. Таким

образом, задача обнаружения дублей сводится к отнесению вектора,

соответствующего анализируемой паре документов, к одному из этих

кластеров. В~\cite{12-sim} для решения этой задачи используется метод

опорных векторов, в~\cite{15-sim} выбирается кластер, центроид которого

находится ближе к вектору.



\section{Предложенный подход к обнаружению дубликатов}



  В целях обеспечения возможности применения разрабатываемого подхода

  к~обработке сообщений из разных предметных областей модель документа

должна предусматривать возможность сравнения сообщений по различным

признакам (см.\ п.~2.2). Окончательное же решение должно приниматься на

основе анализа пары документов с~точки зрения всех критериев. Исходя из

этого, удобно представить пару документов $(d_i,d_j)$ вектором, элементами

которого являются результаты сравнения документов по соответствующим

признакам:

  \begin{equation*}

  \rho_{i,j}=\left( \rho_{i,j}^1,\rho_{i,j}^2, \ldots ,\rho_{i,j}^k\right)\,,

%  \label{e1-sim}

  \end{equation*}

  где $\rho_{i,j}^k$ характеризует сходство документов~$d_i$ и~$d_j$ по

  $k$-му критерию. На основе этого вектора принимается решение о наличии

или отсутствии дублирования. Для этого необходимо задать функцию,

выполняющую интерпретацию~$\rho_{i,j}$, т.\,е.\ определяющую его в~один из

двух классов, один из которых означает наличие дублирования (обозначим его

$M_+$), другой~--- отсутствие ($M_-$):

  \begin{equation}

  D(\rho_{i,j})=\begin{cases}

  1\,,\ &\ \rho_{i,j}\in M_+\,;\\

  0\,, & \ \rho_{i,j}\in M_-\,.

  \end{cases}

  \label{e2-sim}

  \end{equation}



  С учетом выбранного подхода процесс обнаружения дублей можно разбить

на следующие этапы (рис.~2):

  \begin{enumerate}[(1)]

\item построение модели документа, отражающей характеристики новостного

сообщения с~точки зрения каждого из выбранных критериев;

\item сравнение моделей двух документов и~получение вектора~$\rho_{i,j}$;

\item интерпретация вектора с~помощью решающей функции $D(\rho_{i,j})$.

\end{enumerate}



  \begin{figure} %fig2

  \vspace*{1pt}

 \begin{center}

 \mbox{%

 \epsfxsize=127.421mm

 \epsfbox{sim-2.eps}

 }

 \end{center}

 \vspace*{-10pt}

  \Caption{Этапы обнаружения дубликатов}

  \vspace*{-5pt}

  \end{figure}



\section{Модель документа}



  Важным этапом решения задачи является выбор критериев для сравнения

документов. Набор критериев должен быть достаточно выразительным, чтобы

обеспечивать возможность гибкой настройки модели в~соответствии со

спецификой различных предметных областей. Для определения набора

критериев было проведено исследование выборки текстовых документов,

автоматически собираемых из различных ин\-тер\-нет-ис\-точ\-ников.



  В качестве источников были выбраны 35~сайтов по различной тематике:

основные новостные сайты, публикующие материалы

  об\-щест\-вен\-но-по\-ли\-ти\-че\-ской и~экономической тематики, официальные сайты

органов государственной влас\-ти РФ, некоторые сайты органов законодательной

и исполнительной власти субъектов РФ. Такой набор сайтов позволил охватить

значительное число тем и~типов информационных сообщений (ленты новостей,

аналитические статьи и~обзоры, документы правового характера, документы,

содержащие фи\-нан\-со\-во-эко\-но\-ми\-че\-ские показатели).



  Одна тысяча двести пар документов были проанализированы экспертами вручную.

В~результате выполненного анализа были выявлены следующие критерии,

характеризующие модель распознавания дубликатов.



  \vspace*{-2pt}



\subsection{Содержание текста}



  \vspace*{-2pt}



  В большинстве случаев определяющее значение имеет близость текстов.

Преж\-де всего это характерно для общественно-политических новостей, вклад

которых в~суммарную информацию, передаваемую новостным потоком,

определяется оригинальностью их текстового содержания.



  Для обеспечения возможности сравнения составов слов документов модель

включает векторное представление текста сообщения

$$

d_i^w= \left( w_i^1, w_i^2, \ldots , w_i^{N^w}\right)\,,

$$

 где $N^w$~--- общее количество различных

слов во всех документах; $w_i^j$~--- вес $j$-го слова в~$i$-м документе. Если

слово не встречается в~документе, его вес равен нулю. Для остальных слов вес

рассчитывается по методу TF-IDF.



  Векторное представление позволяет использовать для сравнения текстов

прос\-тые алгебраические методы, такие как косинусная мера. Она оценивает

близость векторов на основании значения косинуса угла между ними. При

использовании неотрицательных весов слов косинусная мера принимает

значения в~интервале $[0,\,1]$, поэтому в~качестве оценки различия векторов

используется значение $\rho_{i,j}^w\hm= 1-\mathrm{sim}_{\cos}\left( d_i^w, d_j^w\right)$.



\vspace*{-4pt}



\subsection{Содержание заголовка}



\vspace*{-2pt}



  Отдельно при принятии решения учитываются заголовки, причем их роль

значительно варьируется в~зависимости от предметной области. При

перепечатке политических новостей с~одного сайта на другой нередко

изменяется только заголовок, причем иногда~--- весьма существенно. В~таком

случае наличие дублирования может быть обнаружено на основе близости

остального текста новостей. Однако для сообщений другого типа (например,

правового характера) различие заголовков имеет решающее значение, и~такие

документы не должны признаваться дубликатами, несмотря на близкое

содержание.



  В модели документа заголовок текста~$d_i^t$ представляется аналогично

основному тексту, а~определение близости заголовков~$\rho_{i,j}^t$

выполняется идентично сравнению содержания документов.



\vspace*{-4pt}



\subsection{Предложения и~абзацы}



\vspace*{-2pt}



  Помимо оценки близости текстового содержания документов в~целом

эксперт обращает особое внимание на наличие в~сообщениях идентичных

структурных элементов текстов~--- предложений и~абзацев. Это связано с~тем,

что при перепечатке некоторым источником ранее опубликованного документа

многие из этих элементов переносятся в~текст-дуб\-ли\-кат без изменений.



  Для учета этих критериев документ представляется множеством своих

предложений

$$

d_i^c= \left\{ c_i^1,c_i^2,\ldots ,c_i^{N_i^c}\right\}\,,

$$

 где

$N_i^c$~--- количество предложений в~$i$-м документе. Вес предложения

$\omega(c)$ определяется как сумма весов составляющих его слов.



  При выборе меры близости нужно учесть, что она должна быть

несимметричной: если один из документов получен из другого путем удаления

нескольких предложений, то лишь одно сообщение пары содержит

дополнительную информацию относительно второго, но не наоборот. Для учета

этого требования в~качестве меры сходства была выбрана взвешенная мера

включения:



\vspace*{2pt}



\noindent

  $$

  \mathrm{sim}^c_{\mathrm{inc}} \left( d_i^c,d_j^c\right) = \fr{\sum\limits_{c\in \vert d_i^c\cap

d_j^c\vert}\omega(c)} {\sum\limits_{c\in d_i^c} \omega(c)}\,.

  $$



  \vspace*{-2pt}



  Для оценки различия документов с~точки зрения предложений используются

значения  $\rho_{i,j}^c\hm= 1-\mathrm{sim}^c_{\mathrm{inc}}\left(d_i^c,d_j^c\right)$

и~$\rho_{j,i}^c\hm= 1-\mathrm{sim}^c_{\mathrm{inc}}\left(d_j^c,d_i^c\right)$.



  Построение модели множества абзацев~$d_i^p$ и~определение мер различия

документов с~точки зрения абзацев~$\rho^p_{i,j}$ и~$\rho^p_{j,i}$ выполняется

аналогично.



\vspace*{-4pt}



\subsection{Числовые данные}



\vspace*{-2pt}



  Кроме документов, содержащих только текстовую информацию, часто

встречаются и~те, которые включают числовые данные. В~ряде случаев даже

незначительное изменение этих данных может существенно повлиять на

содержание сообщения. Примером таких документов являются новостные

сообщения из области экономики (новости о состоянии фондового рынка) или

спорта (сообщения о результатах соревнований). Такие документы не должны

признаваться дубликатами даже при полном совпадении их текста.



  Для учета этого признака сообщения представляются множествами чисел,

извлеченных из текста:



\noindent

$$

d_i^n= \left\{ n_i^1, n_i^2, \ldots , n_i^{N_i^{n_d}}\right\}\,,

$$

 где $N_i^{n_d}$~--- количество различных чисел

в~$i$-м документе. Для выполнения оценки сходства таких наборов также

используется мера включения, однако элементы сравниваемых множеств не

являются взвешенными, а~потому учитывается лишь количество одинаковых и~различающихся числовых значений:



\vspace*{2pt}



\noindent

  $$

  \mathrm{sim}^n_{\mathrm{inc}}\left( d_i^n,d_j^n\right) = \fr{\vert d_i^n\cap d_j^n\vert} {\vert

d_i^n\vert}\,.

  $$



  \vspace*{-2pt}



  Отличие одного документа от другого определяется как $\rho_{i,j}^n\hm= 1\hm-

\mathrm{sim}^n_{\mathrm{inc}} \left( d_i^n, d_j^n\right)$.



  Для некоторых предметных областей важен не только состав набора чисел,

но и~порядок их следования в~тексте документа. Это относится,

в~частности, %\linebreak\vspace*{-12pt}

%\pagebreak

%\noindent

 к~спортивным новостям, где разные последовательности одних и~тех же чисел

могут соответствовать различным результатам соревнований (например, два

сета в~теннисном матче, завершившиеся со счетом~6:4 и~4:6). В~таком случае

для представления сообщения используется кортеж $d_i^n\hm= \left( n_i^1,

n_i^2, \ldots , n_i^{N_i^{n_a}}\right)$, где $N_i^{n_a}$~--- общее количество

чисел в~$i$-м документе, а~для сравнения документов~--- расстояние

Да\-ме\-рау--Ле\-вен\-штей\-на $\rho^n_{i,j}\hm= \mathrm{dist}_{\mathrm{DL}}\left( d_i^n, d_j^n\right)$.



\subsection{Фотографии и~ссылки}



  Информация в~сообщении может быть представлена не только текстом или

числовыми данными, но и~различными объектами, включенными в~текст

документа~--- фотографиями, видеороликами, ссылками на сторонние

источники. Присутствие в~документе дополнительных фото- и~видеоматериалов

значительно повышает вероятность его оригинальности, а~наличие ссылок,

напротив, может говорить о заимствовании информации из первоисточников.

%

  Поэтому в~модель включаются компоненты, отражающие количество

присутствующих в~тексте сообщения фотографий $d_i^{\mathrm{im}}$ и~ссылок~$d_i^h$. Различие документов с~точки зрения этих критериев

определяется как разность соответствующих значений:

$$

\rho_{i,j}^{\mathrm{im}}= d_i^{\mathrm{im}}- d_j^{\mathrm{im}}\,;\quad

\rho_{i,j}^h= d_i^h\hm- d_j^h\,.

$$



\subsection{Дата и~время публикации}



  Существенным фактором, влияющим на принятие решения о наличии или

отсутствии дублирования, является разница во времени публикации

сообщений. Так, при дублировании новостных статей перепечатыванию

обычно подвергаются свежие новости, недавно опубликованные на сайте

первоисточника. С~увеличением интервала между моментами появления

документов в~сети вероятность дублирования быстро убывает.



  В модели эта характеристика сообщения представлена посредством

  POSIX-вре\-ме\-ни момента публикации (которое определяется как

количество секунд, прошедших с~полуночи 1~января 1970~г.\ до момента,

когда документ был опубликован источником): $d_i^{\mathrm{dt}}$. Различие между

документами определяется как разность между моментами публикации

сообщений в~секундах:

$$

\rho_{i,j}^{\mathrm{dt}}= d_i^{\mathrm{dt}}- d_j^{\mathrm{dt}}\,.	

$$



\subsection{Авторитетность источника}



  При анализе документов эксперты обращают внимание на источники

сообщений. Статья из авторитетного источника (который обычно публикует

оригинальные материалы) имеет существенно меньшую вероятность быть

признанной дубликатом, чем документ, полученный из источника, регулярно

занимающегося перепечаткой чужих сообщений.



  В модель документа включается компонент, отражающий авторитетность

источника, опубликовавшего это сообщение: $d_i^a\hm\in [0,\,1]$. Это значение

может быть задано экспертом вручную или получено на основе обучающей

выборки как отношение числа оригинальных документов, поступивших от

источника, к~общему числу опубликованных им сообщений.



  Таким образом, модель документа представляет собой совокупность

компонентов, характеризующих сообщение с~точки зрения различных

критериев:

  $$

  d_i= \left( d_i^w, d_i^t, d_i^n, d_i^c, d_i^p, d_i^{\mathrm{im}}, d_i^h, d_i^{\mathrm{dt}},

d_i^a\right)\,.

  $$



\section{Метод обнаружения дубликатов}



\subsection{Интерпретация результата сравнения}



  После получения моделей $d_i$ и~$d_j$ двух документов необходимо

сравнить их и~вынести решение о том, являются ли документы нечеткими

дубликатами. Для этого выполняется анализ схожести сообщений по каждому

из критериев. Результатом сравнения документов с~точки зрения некоторого

критерия~$k$ является значение~$\rho_{i,j}^k$. Выполнив попарно сравнение

компонентов моделей для каждого из критериев, получим вектор

  $$

  \rho_{i,j}=\left( \rho_{i,j}^w,  \rho_{i,j}^t, \rho_{i,j}^n, \rho_{i,j}^c, \rho_{i,j}^p,

  \rho_{i,j}^{\mathrm{im}}, \rho_{i,j}^h, \rho_{i,j}^{\mathrm{dt}}, \rho_{i,j}^a\right)\,.

  $$



  Ввиду вышеуказанной несимметричности мер сходства, используемых для

некоторых критериев, результатом сравнения двух документов являются два

различных вектора~$\rho_{i,j}$ и~$\rho_{j,i}$, характеризующие степень

отличия первого и~второго сообщения друг от друга. Каждый из этих векторов

интерпретируется с~помощью функции $D(\rho)$~(\ref{e2-sim}).



  Для интерпретации результата сравнения необходимо решить задачу

бинарной классификации, т.\,е.\ отнести вектор к классу~$M_+$ или~$M_-$.

Для настройки параметров классификатора используется обучающая

  выборка~--- набор векторов, каждый из которых снабжен меткой $m\hm\in

\{M_+, M_-\}$, обозначающей класс, к которому принадлежит этот вектор.



  Задача бинарной классификации состоит в~том, чтобы для вновь

поступившего на исследование вектора $\rho\hm= \left( \rho^1, \rho^2,\ldots ,

\rho^K\right)$ определить класс, к которому он принадлежит, т.\,е.\

значение~$m$. Для ее решения в~данной работе использовался метод опорных

векторов (SVM~--- support vector machine).

Этот метод основан на построении в~$K$-мер\-ном

пространстве ($K-1$)-мер\-ной гипер\-плос\-кости, разделяющей объекты классов

$M_+$ и~$M_-$. В~зависимости от расположения вектора~$\rho$ относительно

этой гиперплоскости выполняется его отнесение к одному из классов.



  Возможны следующие результаты интерпретации:

  \begin{itemize}

\item $D(\rho_{i,j})=D(\rho_{j,i})=0$. Оба документа признаются оригинальными;

\item $D(\rho_{i,j})\not= D(\rho_{j,i})$. Один из документов является оригиналом, а~второй~--- дублем;

\item $D(\rho_{i,j})=D(\rho_{j,i})=1$. Оба документа являются дублями друг

относительно друга, т.\,е.\ ни один из них не содержит оригинальных данных

относительно другого.

\end{itemize}



  На основе полученных результатов принимается решение о дальнейших

действиях в~отношении документов. Так, в~разработанной системе решалась

задача проверки документов в~момент их поступления от источника, при этом

загружаемые сообщения сравнивались на предмет дублирования с~документами, уже загруженными в~базу. Поэтому интерес представлял лишь

один из результатов интерпретации~--- является ли загружаемый документ

дубликатом ранее полученного сообщения. Однако при обработке готовой

коллекции документов с~целью обнаружения и~устранения дубликатов важно

выявить все пары сообщений, в~которых имеет место дублирование, для чего

требуется использовать оба результата.



\subsection{Метод предварительного отбора кандидатов}



  Предложенный метод выявления нечетких дубликатов подразумевает, что

каждое новое сообщение подвергается сравнению со всеми ранее

загруженными и~при каждом сравнении выполняется расчет близости

документов по множеству критериев. На практике напрямую применять такой

подход не представляется возможным из-за его высокой вычислительной

сложности. Для увеличения быстродействия метода необходимо исключать из

рассмотрения пары новостей, которые слишком сильно отличаются друг от

друга.

%

  С этой целью вышеописанный метод предваряется процедурой отбора пар

документов, являющихся потенциальными дубликатами:

  $$

  F(d_i,d_j)=\begin{cases}

  1, &\ \mbox{если документы~--- потенциальные дубликаты};\\

  0, &\ \mbox{если документы не являются дубликатами}.

  \end{cases}

  $$



  Пары, признанные процедурой кандидатами в~дубликаты, должны

проверяться основным методом, остальные же сообщения далее не

рассматриваются (рис.~3).



  \begin{figure} %fig3

  \vspace*{1pt}

 \begin{center}

 \mbox{%

 \epsfxsize=126.897mm

 \epsfbox{sim-3.eps}

 }

 \end{center}

 \vspace*{-10pt}

  \Caption{Двухэтапный подход к обнаружению дубликатов}

  \vspace*{-2pt}

  \end{figure}



  Эта процедура, по сути, также решает задачу обнаружения дубликатов,

причем основными требованиями, предъявляемыми к ней, являются

минимальная вычислительная сложность, максимальная полнота, а~также

высокая эффективность (под эффективностью понимается отношение числа пар

документов, отброшенных на этапе предварительного отбора, к исходному

числу пар).



  Ввиду наличия множеств взвешенных слов документов было решено

положить в~основу процедуры сравнение наборов наиболее значимых слов:

$s\hm= \left\{ s^1, s^2, \ldots  , s^A\right\}$, где $A$~--- мощность сравниваемых

наборов, определяемая опытным путем. При этом документы признаются

кандидатами в~дубликаты, если мощность пересечения их наборов не меньше

заданного порогового значения:



\vspace*{2pt}



\noindent

  $$

  F\left(d_i,d_j\right) =\begin{cases}

  1, &\ \vert s_i\cap s_j\vert \geq T\,;\\

  0, &\ \vert s_i\cap s_j \vert <T\,.

  \end{cases}

  $$



  Мощность набора $A$ и~пороговое значение $T\hm\in [0;\,A]$ являются

настраиваемыми параметрами, выбираемыми исходя из желаемых показателей

полноты и~эффективности.





\section{Экспериментальная проверка метода}



  В рамках данной работы были проведены эксперименты, направленные на

анализ качества работы разработанной подсистемы обнаружения дубликатов,

реализующей предложенный метод. Все эксперименты проводились на ПЭВМ

со следующими основными параметрами: процессор Intel Core 2 Duo 2,2 ГГц,

объем ОЗУ~--- 2 ГБ.



  Целью первого эксперимента была оценка качества метода предварительного

отбора потенциальных дубликатов в~зависимости от параметров~$A$ и~$T$ на

примере анализа общественно-политических новостей. Тестирование

проводилось в~течение суток. В~качестве входных данных использовались

1\,427~документов, извлеченных с~20~новостных сайтов. Каждый из

документов подвергался сравнению с~17\,835~загруженными ранее

сообщениями. В~общей сложности было обработано 25\,450\,545~пар

документов, из которых~676~были дубликатами.



  В таблице приведены результаты работы метода предварительного отбора

при использовании трех конфигураций параметров процедуры.



  \begin{table}\small

  \begin{center}

  \vspace*{2ex}



  \tabcolsep=5pt

\begin{tabular}{|c|c|c|c|c|c|c|c|c|}

\multicolumn{9}{c}{Оценка качества метода отбора кандидатов}\\

\multicolumn{9}{c}{\ }\\[-6pt]

  \hline

  \tabcolsep=0pt\begin{tabular}{c}Конфи-\\ гурация\end{tabular} &

  \tabcolsep=0pt\begin{tabular}{c}

$A$ \end{tabular} &

\tabcolsep=0pt\begin{tabular}{c}$T$\end{tabular}&

\tabcolsep=0pt\begin{tabular}{c}Распо-\\ знанные\\ дубли-\\ каты\end{tabular} &

\tabcolsep=0pt\begin{tabular}{c}Ложные\\ срабаты-\\ вания\end{tabular}&

\tabcolsep=0pt\begin{tabular}{c}Пропу-\\ щенные\\ дубли-\\ каты\end{tabular}&

 Полнота& Точность&

 \tabcolsep=0pt\begin{tabular}{c}Эффектив-\\ ность\end{tabular}\\

\hline

&&&&&&&&\\[-9pt]

1& 2& 2&352& \hphantom{9}482  & 324& 52,1\% & 42,2\%& 99,99\%\\

2& 6&2 &625& 4231 & \hphantom{9}51 & 92,4\% & 12,8\%& 99,98\%\\

3 & 8&2  &651& 8724 & \hphantom{9}25 & 96,3\% & \hphantom{9}6,9\% & 99,96\%\\

\hline

\end{tabular}

\end{center}

\end{table}



  Использование более мягких критериев фильтрации (конфигурация~3)

предсказуемо приводит к увеличению полноты распознаваемых дубликатов и~к

уменьшению эффективности процедуры по сравнению с~жесткими

настройками (конфигурация~1). Исходя из требований к процедуре,

оптимальной была признана конфигурация~3: сравнению должны

подвергаться~8~наиболее значимых слов документа, при

совпадении~2~и~более документы считаются кандидатами в~дубликаты. Такие

параметры позволяют получить высокую полноту, равную 96,3\%, при этом

дальнейшая проверка основным методом требуется лишь для 0,04\% от общего

числа пар.



  При таких настойках дубликаты составляют лишь 6,9\% пар, прошедших

отбор. Столь низкая точность процедуры доказывает необходимость

дополнительного анализа отобранных пар с~помощью основного метода.



  В рамках второго эксперимента выполнялась оценка качества основного

метода. Для тестирования использовались 26\,036~документов, извлеченных с~20~новостных сайтов. С помощью метода фильтрации было отобрано

2650~пар-кан\-ди\-да\-тов, каждая из которых была проанализирована

экспертами на предмет наличия дублирования. Часть пар использовалась для

обучения, на остальных выполнялось тестирование метода. Целью

эксперимента было определение зависимости показателей качества (точности,

полноты и~$F$-ме\-ры) от мощности обучающей выборки и~от учитываемых

критериев. Полученные зависимости представлены на рис.~4.







  Как видно из рис.~4, при использовании 400 обучающих примеров

происходит насыщение и~с дальнейшим увеличением обучающей выборки

качество работы метода не улучшается. Таким образом, для обучения системы

достаточно 400 пар документов, размеченных экспертами.



  Проведенный эксперимент доказывает целесообразность

многокритериального сравнения документов: при использовании всех

критериев достигаются более высокие показатели качества ($F$-ме\-ра в~зоне

насыщения равна 0,82), чем при анализе документов только с~точки зрения слов

(0,67) или абзацев (0,64).



  \begin{figure} %fig4

  \vspace*{1pt}

 \begin{center}

 \mbox{%

 \epsfxsize=122.581mm

 \epsfbox{sim-4.eps}

 }

 \end{center}

 \vspace*{-10pt}

  \Caption{Зависимость точности (\textit{1}), полноты (\textit{2})

  и~$F$-меры~(\textit{3}) от мощности обучающей выборки при использовании

различных наборов критериев: (\textit{а})~использование близости составов слов;

(\textit{б})~использование схожести абзацев;

(\textit{в})~использование всех критериев}

\vspace*{-6pt}

  \end{figure}



  При более подробном анализе примеров, на которых метод совершает

ошибку, оказалось, что в~большинстве случаев (76\%) ошибка связана с~человеческим фактором: каждый эксперт, участвовавший в~подготовке

обучающей и~тестовой выборок, имеет свое представление о том, какие

документы являются информативно необходимыми, а~какие~--- нет, что

приводит к конфликтам в~суждениях экспертов. Лишь в~24\% случаев

ошибочность решения, принятого методом, не вызывала сомнений.



  Проведенный эксперимент также показал, что среднее время, затрачиваемое

на анализ пары документов основным методом, составляет 5~мс. С~учетом

высокой эффективности процедуры фильтрации это означает, что система

способна обрабатывать 10\,000~документов в~час при наличии в~базе

200\,000~ранее загруженных сообщений, с~которыми требуется сравнивать

новые документы.



\section{Направления дальнейших исследований}



  Помимо обнаружения и~устранения дубликатов, предлагаемый метод может

быть использован для решения других задач интеллектуального анализа

текстов. При соответствующей настройке набора учитываемых критериев и~порога

бли\-зости документов, необходимого для вынесения решения о наличии

дублирования, метод может быть применен для решения общей задачи

обнаружения документов, близ\-ких по содержанию к заданному. Это позволит,

в частности, выполнять формирование подборок тематически близких

документов, а~также сообщений, которые с~большой долей вероятности связаны

с~общим событием.



  Еще одним перспективным направлением развития метода является

снабжение его возможностью не только обнаружения наличия или отсутствия

дублирования, но и~выделения в~тексте близких по содержанию сообщений

фрагментов с~оригинальной (недублированной) информацией.



\section{Заключение}



  В работе предложен метод обнаружения и~устранения нечеткого

дублирования в~потоке текстовых сообщений. В~его основе лежит отнесение

пар документов к~классу дубликатов или недубликатов с~помощью метода

опорных векторов.



  Предлагаемый метод обладает высокой гибкостью благодаря возможности

его настройки для обработки сообщений из различных предметных областей.

Это достигается посредством включения в~модель документа компонентов,

отражающих критерии, которыми руководствуются эксперты при анализе

текстовых коллекций вручную.



  Для обеспечения низкой вычислительной сложности предложена процедура

отбора пар-кан\-ди\-да\-тов на основе сравнения наборов наиболее значимых

термов документов. Это позволяет применять основной метод лишь к~парам

документов, прошедшим отбор.



  Представленный метод был апробирован при решении задачи анализа потока

текстовых сообщений, загружаемых из открытых ин\-тер\-нет-ис\-точ\-ни\-ков,

с~целью устранения документов, являющихся дубликатами ранее загруженных

материалов.



{\small\frenchspacing

{%\baselineskip=10.8pt

\addcontentsline{toc}{section}{Литература}

\begin{thebibliography}{99}

\bibitem{1-sim}

\Au{Ландэ Д.\,В., Дармохвал А.\,Т., Морозов А.\,Ю.} Подход к выявлению дублирования

сообщений в~новостных информационных потоках~/\!\!/ Электронные библиотеки:

перспективные методы и~технологии, электронные коллекции: Тр. 8-й Всеросс. науч. конф.

(RCDL'2006).~---  Ярославль: ЯрГУ им. П.\,Г.~Демидова, 2006. C.~115--119.

\bibitem{2-sim}

\Au{Зеленков Ю.\,Г., Сегалович И.\,В.} Сравнительный анализ методов определения нечетких

дубликатов для Web-документов~/\!\!/ Электронные библиотеки: перспективные методы и

технологии, электронные коллекции: Тр. 9-й Всеросс. науч. конф. (RCDL'2007).~---

Пе\-ре\-славль-За\-лес\-ский: Университет города Переславля, 2007. С.~166--174.

\bibitem{3-sim}

\Au{Андреев А.\,М., Березкин Д.\,В., Козлов И.\,А., Симаков~К.\,В.} Метод обнаружения

изменений структуры веб-сайтов в~системе сбора новостной информации~/\!\!/ Электронные

библиотеки: перспективные методы и~технологии, электронные коллекции: Тр. 14-й Всеросс.

науч. конф. (RCDL'2012).~--- Пе\-ре\-славль-За\-лес\-ский: Университет города Переславля,

2012. С.~124--133.

\bibitem{4-sim}

\Au{Блох М.\,Я.} Теоретические основы грамматики: Учебник.~--- 2-е изд., испр.~--- М.: Высшая

школа, 2000. 160~с.

\bibitem{5-sim}

\Au{Chowdhury A., Frieder O., Grossman D., McCabe~M.} Collection statistics for fast duplicate

document detection~/\!\!/ ACM Trans. Inform. Syst. (TOIS), 2002. Vol.~20. No.\,2. P.~171--191.

\bibitem{6-sim}

\Au{Ilyinsky S., Kuzmin M., Melkov A., Segalovich~I.} An efficient method to detect duplicates of Web

documents with the use of inverted index~/\!\!/ 11th World Wide Web Conference (International)

(WWW'2002) Proceedings.~--- Honolulu, Hawaii, USA: ACM Press, 2002. {\sf

http://www2002.org/CDROM/poster/187/index.html}.

\bibitem{7-sim}

\Au{Brin S., Davis J., Garcia-Molina H.} Copy detection mechanisms for digital documents~/\!\!/

ACM SIGMOD Record, 1995. Vol.~24. No.\,2. P.~398--409.

\bibitem{8-sim}

\Au{Broder A.} Algorithms for duplicate documents. {\sf

http://www.cs.princeton.edu/\linebreak courses/archive/spr05/cos598E/bib/Princeton.pdf}.

\bibitem{9-sim}

\Au{Manber U.} Finding similar files in a large file system~/\!\!/  USENIX Winter 1994 Technical

Conference Proceedings.~--- Berkeley, CA, USA: USENIX Association, 1994. P.~1--10.

\bibitem{10-sim}

\Au{Shivakumar N., Garcia-Molina H.} SCAM: A~copy detection mechanism for digital

documents~/\!\!/ 2nd Conference (International) in Theory and Practice of Digital Libraries (DL'95)

Proceedings.~--- Austin, TX, USA, 1995. {\sf http://ilpubs.stanford. edu:8090/95/1/1995-28.pdf}.

\bibitem{11-sim}

\Au{Шевелёв О.\,Г.} Методы автоматической классификации текстов

на естественном языке.~--- Томск: ТМЛ-Пресс, 2007. 144~с.



\bibitem{13-sim} %12

\Au{Ko{\hspace*{-1mm}\ptb{\l}}cz A., Chowdhury A., Alspector J.} Improved robustness of signature-based near-

replica detection via lexicon randomization~/\!\!/ 10th ACM SIGKDD Conference (International) on

Knowledge Discovery and Data Mining Proceedings.~--- ACM, 2004. P.~605--610.

\bibitem{14-sim} %13

\Au{Fetterly D., Manasse M., Najork M., Wiener J.} A large-scale study of the evolution of web

pages~/\!\!/ 12th Conference (International) on World Wide Web (WWW'03) Proceedings.~--- New

York, NY, USA: ACM Press, 2003. P.~669--678.



\bibitem{12-sim} %14

\Au{Bilenko M., Mooney R.\,J.} Adaptive duplicate detection using learnable string similarity

measures~/\!\!/ 9th ACM SIGKDD Conference (International) on Knowledge Discovery and Data

Mining Proceedings.~--- ACM, 2003. P.~39--48.

\bibitem{15-sim}

\Au{Князева А.\,А., Турчановский И.\,Ю., Колобов О.\,С.} Выявление дубликатов

в~биб\-лио\-графических базах данных~/\!\!/ Электронные библиотеки: перспективные методы и

технологии, электронные коллекции: Тр. 15-й Всеросс. науч. конф. (RCDL'2013).~--- Ярославль:

ЯрГУ им. П.\,Г.~Демидова, 2013. С.~207--213.

\end{thebibliography}

} }





%\vspace*{-3pt}



\hfill{\small\textit{Поступила в~редакцию 30.12.14}}









%\vspace*{9pt}



%\hrule



%\vspace*{2pt}



%\hrule







\newpage



\vspace*{-24pt}





\def\tit{MULTICRITERIA METHOD FOR DETECTING NEAR-DUPLICATES

IN~A~STREAM OF~TEXT MESSAGES}



\def\titkol{Multicriteria method for detecting near-duplicates

in~a~stream of~text messages}





\def\aut{A.~Andreev, D.~Berezkin, I.~Kozlov, and~K.~Simakov}

\def\autkol{A.~Andreev, D.~Berezkin, I.~Kozlov, and~K.~Simakov}





\titel{\tit}{\aut}{\autkol}{\titkol}



\vspace*{-6pt}



\noindent

Bauman Moscow State Technical University, 5~Baumanskaya 2nd Str.,Moscow 105005,

Russian Federation







\def\leftfootline{\small{\textbf{\thepage}

\hfill Sistemy i Sredstva Informatiki~---

Systems and Means of Informatics 2015 vol~25 no\,1}

}%

 \def\rightfootline{\small{Sistemy i Sredstva Informatiki~---

 Systems and Means of Informatics   2015   vol~25   no\,1

\hfill \textbf{\thepage}}}





\Abste{The problem of near-duplicate detection in a stream of text messages is

considered. A model of a text document and a~multicriteria duplicate identification

method is proposed. The model provides flexible adjustment for different domains.

The method is based on binary classification using support vector machine. The paper

also provides a method of candidates prefiltration in order to ensure high efficiency

of the approach. Several experiments with data obtained from a stream of news

articles were carried out. The results show feasibility of the suggested approach.}



\KWE{near-duplicate detection; similarity measure; binary classification}



\DOI{10.14357\!/\!08696527150103}



%\vspace*{-24pt}



%\Ack

%\noindent

%Статья рекомендована к публикации в~журнале Программным комитетом конференции <<Электронные

%библиотеки: перспективные методы и~технологии, электронные коллекции>> (``Digital Libraries: Advanced

%Methods and Technologies,'' RCDL-2014), 13--16~октября 2014~года, г.~Дубна, РФ.





\renewcommand{\bibname}{\protect\rmfamily References}

%\renewcommand{\bibname}{\large\protect\rm References}



{\small\frenchspacing

{%\baselineskip=10.8pt

\addcontentsline{toc}{section}{References}

\begin{thebibliography}{99}

\bibitem{1-sim-1}

\Aue{Lande, D.\,V., A.\,T. Darmokhval, and A.\,Yu.~Morozov}. 2006. Podkhod k

vyyavleniyu dublirovaniya soobshcheniy v novostnykh informatsionnykh potokakh [The

approach to duplication detection in news information flows].

\textit{Elektronnye biblioteki: Perspektivnye metody i~tekhnologii, elektronnye

kollektsii:

Tr. 8-y Vseross. konf. (RCDL2006)} [Digital Libraries: Advanced Methods and

Technologies, Digital Collections: 8th All-Russian Scientific Conference RCDL-2006

Proceedings].  Suzdal. 115--119.

\bibitem{2-sim-1}

\Aue{Zelenkov,~Yu.\,G., and I.\,V.~Segalovich}. 2007. Sravnitel'nyy analiz metodov

opredeleniya nechetkikh dublikatov dlya Web-dokumentov [Comparative analysis of

near-duplicate detection methods of Web documents]. \textit{Elektronnye biblioteki:

Perspektivnye metody i~tekhnologii, elektronnye kollektsii: Tr. 9-y Vseross. nauch.

konf.  (RCDL'2007)} [Digital Libraries: Advanced Methods and Technologies, Digital

Collections: 9th All-Russian Scientific Conference RCDL'2007 Proceedings].

Pereslavl'-Zalessky. 166--174.

\bibitem{3-sim-1}

\Aue{Andreev, A.\,M., D.\,V.~Berezkin, I.\,A.~Kozlov, and K.\,V.~Simakov.} 2012.

Metod obnaruzheniya izmeneniy struktury veb-saytov v~sisteme sbora novostnoy

informatsii [The method of detecting structure changes of news web sites].

\textit{Elektronnye biblioteki: Perspektivnye metody i~tekhnologii, elektronnye

kollektsii: Tr. 14-y Vseross. nauch. konf. (RCDL2012)} [Digital Libraries: Advanced

Methods and Technologies, Digital Collections:  14th All-Russian Scientific Conference

RCDL2012 Proceedings]. Pereslavl'-Zalessky. 124--133.

\bibitem{4-sim-1}

\Aue{Blokh, M.\,Ya.} 2000. \textit{Teoreticheskie osnovy grammatiki}

[Theoretical Foundations of Grammar]. 2nd ed. Moscow: Vysshaya Shkola.

160~p.

\bibitem{5-sim-1}

\Aue{Chowdhury, A., O.~Frieder, D.~Grossman, and M.~McCabe.} 2002. Collection statistics for fast duplicate

document detection. \textit{ACM Trans. Inform. Syst. (TOIS)} 20(2):171--191.

\bibitem{6-sim-1}

\Aue{Ilyinsky, S., M.~Kuzmin, A.~Melkov, and I.~Segalovich}. 2002. An efficient method to detect duplicates of Web

documents with the use of inverted index. \textit{11th World Wide Web Conference

(International) (WWW'2002) Proceedings}. Honolulu, Hawaii: ACM Press.

Available at: {\sf

http://www2002.org/CDROM/poster/187/index.html} (accessed April~1, 2015).

\bibitem{7-sim-1}

\Aue{Brin, S., J.~Davis, and H.~Garcia-Molina}. 1995. Copy detection mechanisms for

digital documents. \textit{ACM SIGMOD Record} 24(2):398--409.

\bibitem{8-sim-1}

\Aue{Broder, A.} Algorithms for duplicate documents. Available at:  {\sf

http://www.cs.princeton. edu/courses/archive/spr05/cos598E/bib/Princeton.pdf}

(accessed December~17, 2014).

\bibitem{9-sim-1}

\Aue{Manber, U.} 1994. Finding similar files in a large file system.

\textit{USENIX Winter 1994 Technical

Conference Proceedings}. Berkeley, CA, USA: USENIX Association. 1--10.

\bibitem{10-sim-1}

\Aue{Shivakumar, N., and H. Garcia-Molina}. 1995. SCAM: A~copy detection

mechanism for digital documents. \textit{2nd Conference in Theory and Practice of

Digital Libraries (DL'95) Proceedings.} Available at: {\sf

http://ilpubs.stanford.edu:8090/95/1/1995-28.pdf} (accessed April~1, 2015).



\bibitem{11-sim-1}

\Aue{Shevelev, O.\,G.} 2007. \textit{Metody avtomaticheskoy klassifikatsii tekstov

na estestvennom yazyke} [Methods of automatic classification of

natural language texts]. Tomsk: TML-Press. 144~p.





\bibitem{13-sim-1} %12

\Aue{Ko{\hspace*{-1mm}\ptb{\l}}cz,~A., A.~Chowdhury, and J.~Alspector}. 2004. Improved

robustness of signature-based near-replica detection via lexicon randomization.

\textit{10th ACM SIGKDD Conference (International) on Knowledge Discovery and

Data Mining Proceedings}. 605--610.

\bibitem{14-sim-1} %13

\Aue{Fetterly, D., M.~Manasse, M.~Najork, and J.~Wiener}. 2003. A~large-scale study of the evolution of web

pages. \textit{12th Conference (International) on World Wide Web (WWW'03)

Proceedings}. New York, NY: ACN Press. 669--678.



\bibitem{12-sim-1} %14

\Aue{Bilenko, M., and R.\,J. Mooney}. 2003. Adaptive duplicate detection using

learnable string similarity measures. \textit{9th ACM SIGKDD Conference

(International) on Knowledge Discovery and Data Mining Proceedings}. 39--48.



\bibitem{15-sim-1} %15

\Aue{Knyazeva, A.\,A., I.\,Yu. Turchanovskiy, and O.\,S.~Kolobov}. 2013.

Vyyavlenie dublikatov v bibliograficheskikh bazakh dannykh [Identification of

duplicates in bibliographic databases]. \textit{Elektronnye biblioteki: Perspektivnye

metody i~tekhnologii, elektronnye kollektsii: Tr. 15-y Vseross. nauch. konf.

(RCDL2013)} [Digital Libraries: Advanced Methods and Technologies, Digital

Collections: 15th All-Russian Scientific Conference RCDL2013 Proceedings].

Yaroslavl. 207--213.

\end{thebibliography}

} }





\vspace*{-6pt}



\hfill{\small\textit{Received December 30, 2014}}





  \Contr



\noindent

\textbf{Andreev Ark M.} (b.\ 1943)~--- Candidate of Science (PhD) in technology,

Associate professor, Bauman Moscow State Technical University,

5~Baumanskaya 2nd Str., Moscow 105005, Russian Federation; arkandreev@gmail.com



\pagebreak



\noindent

\textbf{Berezkin Dmitry V.} (b.\ 1966)~---

Candidate of Science (PhD) in technology, senior scientist, Bauman Moscow State

Technical University, 5~Baumanskaya 2nd Str., Moscow 105005, Russian Federation;

dmitryb2007@yandex.ru



\vspace*{3pt}



\noindent

\textbf{Kozlov Ilya A.} (b.\ 1989)~---

PhD student, Bauman Moscow State Technical University, 5~Baumanskaya 2nd

Str., Moscow 105005, Russian Federation; \mbox{kozlovilya89@gmail.com}



\vspace*{3pt}



\noindent

\textbf{Simakov Konstantin V.} (b.\ 1980)~---

Candidate of Science (PhD) in technology, senior scientist, Bauman Moscow State

Technical University, 5~Baumanskaya 2nd Str., Moscow 105005, Russian Federation;

skv@ixlab.ru



 \label{end\stat}





\renewcommand{\bibname}{\protect\rm Литература}

